% ==============================================================================
% CHAPTER 2: LITERATURE REVIEW & PROBLEM FORMULATION
% ==============================================================================
\chapter{Literature Review and Problem Formulation}
\label{chap:literature_review}

\section{Review of Existing Literature}
\label{sec:lit_review}

Sentiment analysis, also termed opinion mining, represents one of the most vigorously investigated domains in Natural Language Processing (NLP) and computational linguistics. Over the past two decades, the discipline has evolved through three distinct algorithmic paradigms: (1) \textit{Lexicon-based and Rule-based Heuristics}, (2) \textit{Classical Statistical Machine Learning}, and (3) \textit{Deep Bidirectional Pre-trained Transformers and Large Language Models}. This section conducts a rigorous, multi-faceted survey of foundational and contemporary literature across these paradigms, paying particular attention to their application within computer-mediated communication, online gaming communities, and competitive esports.

\subsection{Foundational Statistical NLP and Classical Machine Learning}

The foundational foundations of computational sentiment classification were established by \textbf{Pang, Lee, and Vaithyanathan (2002)} in their seminal paper evaluating machine learning techniques for movie review categorization \cite{pang2002thumbs}. Their investigation benchmarked three standard supervised classifiers: Naive Bayes, Maximum Entropy classification, and Support Vector Machines (SVM). Pang et al.\ demonstrated that while traditional topic classification tasks benefit from simple bag-of-words boolean presence, sentiment polarity is governed by nuanced contextual modifiers and rhetorical negations. Their findings showed that Linear SVMs using unigram feature vectors achieved the highest classification accuracy (82.9\%), outperforming Naive Bayes (81.5\%) and human-crafted keyword lists (69.0\%). However, the authors explicitly noted that classical statistical models suffer severe degradation when confronted with sarcastic humor, complex discourse structures, and domain-dependent sentiment inversion.

The mathematical formulation of Support Vector Machines in text categorization was formalized by \textbf{Joachims (1998)} \cite{joachims1998text}. Joachims proved that text classification problems exhibit high-dimensional feature spaces (frequently exceeding $10^4$ unique n-gram features) where nearly all features carry non-zero informational value. Furthermore, text vectors are inherently sparse, meaning that individual document representations contain very few active non-zero elements. Joachims demonstrated that linear hyperplanes possess strong theoretical protection against over-fitting in high dimensions via maximum-margin separation, establishing Linear Support Vector Classifiers (LinearSVC) as the canonical, gold-standard classical baseline for computational text categorization.

In parallel, \textbf{McCallum and Nigam (1998)} conducted a foundational comparison between the multi-variate Bernoulli model and the multinomial model for Naive Bayes text classification \cite{mccallum1998comparison}. Their empirical evaluations proved that the Multinomial Naive Bayes (MNB) formulation—which accounts for token frequencies rather than mere binary occurrence—consistently outperforms the Bernoulli model on realistic document lengths. By incorporating Laplace add-one smoothing, MNB avoids zero-probability collapse when encountering rare or previously unseen vocabulary, delivering exceptional inference speeds suitable for microsecond-scale execution.

Addressing the challenge of informal microblogging text, \textbf{Go, Bhayani, and Huang (2009)} pioneered distant supervision techniques for Twitter sentiment classification \cite{go2009twitter}. By automatically collecting training corpora based on emoticons (\textit{e.g.}, ``:)'' as positive and ``:('' as negative), they eliminated the requirement for expensive manual annotation. Evaluating Naive Bayes, Maximum Entropy, and SVM over 1.6 million tweets, they achieved 82.2\% accuracy using SVM with unigram and bigram features. Nonetheless, their study revealed that distant supervision labels introduce considerable label noise, particularly when emoticons are deployed ironically or sarcastically.

\subsection{Lexicon-Based and Rule-Augmented Sentiment Analysis}

Prior to the dominance of deep learning, lexicon-based approaches provided transparent, unsupervised sentiment scoring. \textbf{Hutto and Gilbert (2014)} introduced \textbf{VADER} (Valence Aware Dictionary and sEntiment Reasoner), a rule-based model specifically engineered for social media discourse \cite{hutto2014vader}. VADER combined a curated gold-standard lexical dictionary with five explicit heuristic rules governing:
\begin{enumerate}
    \item Punctuation amplification (\textit{e.g.}, exclamation points);
    \item Capitalization emphasis (\textit{e.g.}, ``GREAT'' vs. ``great'');
    \item Degree modifier shifts (\textit{e.g.}, ``barely good'' vs. ``extremely good'');
    \item Contrastive conjunction flips (\textit{e.g.}, ``The graphics are great, but the balance is terrible'');
    \item Tri-gram negation handling (\textit{e.g.}, ``not really smart'').
\end{enumerate}
While VADER achieves high efficiency on Twitter and product reviews, \textbf{Mohammad and Turney (2013)} demonstrated in their work on crowdsourced emotion lexicons that static lexicons fail to generalize to specialized subcultures \cite{mohammad2013crowdsourcing}. When evaluated on video game discourse, VADER regularly misinterprets gaming vernacular, treating words like ``crazy'', ``sick'', and ``killing'' as severe negative sentiment, while failing to recognize toxic acronyms or title-specific terms.

The international benchmarking of sentiment analysis on computer-mediated communication was formalized through the annual \textbf{SemEval} competitions, notably documented by \textbf{Rosenthal, Farra, and Nakov (2017)} in \textit{SemEval-2017 Task 4: Sentiment Analysis in Twitter} \cite{rosenthal2017semeval}. Over five iterations, the SemEval benchmark tracked the performance transition from hand-engineered feature sets to deep neural networks, proving that capturing semantic context across non-standard spelling and emoji sequences requires representation learning rather than static lexical matching.

\subsection{Deep Learning and Transformer Architectures}

The breakthrough in modern contextual NLP was catalyzed by \textbf{Vaswani et al. (2017)} through the introduction of the \textbf{Transformer} architecture in their landmark paper, \textit{``Attention Is All You Need''} \cite{vaswani2017attention}. By discarding recurrence and convolutions in favor of multi-head scaled dot-product self-attention mechanisms, the Transformer enabled full parallelization of sequence processing and captured bidirectional token dependencies across arbitrarily long textual spans.

Building upon this architecture, \textbf{Devlin et al. (2019)} developed \textbf{BERT} (Bidirectional Encoder Representations from Transformers) \cite{devlin2019bert}. Pre-trained on massive unlabeled corpora (BooksCorpus and English Wikipedia) via Masked Language Modeling (MLM) and Next Sentence Prediction (NSP), BERT achieved state-of-the-art results across eleven NLP benchmarks, including GLUE, MultiNLI, and SQuAD. BERT established that pre-trained contextual representations, when fine-tuned on task-specific classification heads, drastically outperform models trained from scratch.

However, the computational footprint of BERT-Base (110 million parameters, 440\,MB disk footprint, 200\,ms CPU inference per batch) poses formidable operational hurdles for real-time web deployment. To mitigate this computational bottleneck, \textbf{Sanh et al. (2019)} introduced \textbf{DistilBERT}, a distilled version of BERT that compresses the model by 40\% while retaining 97\% of its language understanding capabilities \cite{sanh2019distilbert}. DistilBERT employs knowledge distillation during pre-training, where a compact student network (6 transformer layers instead of 12) is trained under the supervision of a larger teacher network using a triple loss function combining cross-entropy, distillation loss, and cosine embedding loss. Sanh et al.\ proved that DistilBERT runs 60\% faster than BERT while requiring significantly less memory, making it the ideal deep learning backbone for low-latency web applications.

Further refinements in transformer training methodology were demonstrated by \textbf{Liu et al. (2019)} in \textbf{RoBERTa} (Robustly Optimized BERT Approach) \cite{liu2019roberta}. Liu et al.\ proved that BERT was severely undertrained, showing that removing the NSP task, training for longer with larger mini-batches, and dynamically altering masking patterns substantially improves downstream classification performance.

\subsection{Esports Analytics and Player Behavior Mining}

While general sentiment analysis is thoroughly documented, academic literature specifically addressing natural language sentiment and community dynamics within competitive esports has emerged only recently.

\textbf{Mora-Cantallops and Sicilia (2018)} published a comprehensive investigation into player community sentiment and toxic behavior in MOBA titles, focusing on \textit{League of Legends} \cite{mora2018exploring}. Their empirical study examined tens of thousands of official forum discussions, discovering a strong statistical correlation between perceived patch imbalance, champion win-rate variance, and toxic community discourse. They emphasized that player frustration is heavily localized around specific playable characters, arguing that game designers must monitor entity-level sentiment rather than aggregate community sentiment.

In the realm of predictive esports modeling, \textbf{Makarov et al. (2018)} explored competitive match outcome forecasting in \textit{Dota 2} and \textit{CS:GO} using in-game telemetry and team composition data \cite{makarov2018predicting}. Their work proved that pre-match hero drafting and player historical performance metrics account for over 65\% of match outcome variance. However, Makarov et al.\ explicitly identified the absence of community and team morale sentiment as a primary limitation in their predictive pipelines.

Addressing toxicity and communication in team shooters, \textbf{Hodge et al. (2019)} conducted an empirical survey of toxic behavior and communication dynamics in competitive first-person shooters \cite{hodge2019toxicity}. Their findings revealed that tactical voice and text communications in competitive FPS titles (\textit{Counter-Strike: Global Offensive} and \textit{Overwatch}) follow distinct hierarchical patterns: strategic utility coordination is dominant in early rounds, but rapidly devolves into aggressive, accusatory blame-shifting (toxic ``diffing'') when defensive holds fail. Hodge et al.\ concluded that automated toxicity and sentiment detection tools must understand tactical gaming jargon to distinguish genuine tactical advice from hostile harassment.

More recently, \textbf{Schubert, Drachen, and Mahlmann (2020)} explored esports analytics pipelines, advocating for integrated architectures that unify quantitative player telemetries with qualitative community discourse mining \cite{schubert2020esports}. They noted that existing esports organizations suffer from fragmented analytical toolsets, where match performance trackers operate in complete isolation from player sentiment streams.

\subsection{Explainable AI and Multi-Tier Inference Cascades}

As deep neural networks and transformer models have become ubiquitous, the requirement for Explainable Artificial Intelligence (XAI) and model interpretability has become paramount. \textbf{Ribeiro, Singh, and Guestrin (2016)} introduced \textbf{LIME} (Local Interpretable Model-agnostic Explanations), which explains the predictions of any complex machine learning classifier by learning an interpretable linear surrogate model locally around the prediction \cite{ribeiro2016why}. Similarly, \textbf{Lundberg and Lee (2017)} introduced \textbf{SHAP} (SHapley Additive exPlanations), unifying cooperative game-theoretic Shapley values with local feature attribution methods \cite{lundberg2017unified}. 

In the domain of efficient production inference, \textbf{Sun et al. (2020)} investigated multi-tier model cascading and adaptive routing in transformer deployment \cite{sun2020mobilebert}. Their research demonstrated that cascading inference—wherein lightweight statistical classifiers process the majority of high-confidence standard queries while deep neural models or LLMs are invoked exclusively for low-confidence or high-ambiguity edge cases—reduces total system compute costs by over 70\% while maintaining near-transformer accuracy bounds across the entire traffic volume.

\section{Summary of Literature Survey}
\label{sec:lit_table}

To synthesize the scholarly foundation of this project, Table~\ref{tab:literature_survey} summarizes fifteen seminal academic works across classical NLP, deep contextual models, esports analytics, and explainable AI, highlighting their core technical contributions and primary limitations.

\begin{table}[htbp]
\centering
\scriptsize
\caption{Comprehensive literature survey of foundational and modern NLP, esports analytics, and XAI studies.}
\label{tab:literature_survey}
\begin{tabularx}{\textwidth}{P{2.6cm} P{2.8cm} L L}
\toprule
\textbf{Author \& Year} & \textbf{Domain / Architecture} & \textbf{Core Contribution \& Findings} & \textbf{Technical Limitations} \\
\midrule
Pang et al. (2002) \cite{pang2002thumbs} & Classical ML (SVM, NB, MaxEnt) & First comprehensive benchmark for sentiment analysis; proved Linear SVM achieves 82.9\% on movie reviews. & Failed on sarcasm, long-distance negation, and domain slang. \\
\addlinespace
Joachims (1998) \cite{joachims1998text} & Support Vector Machines & Proved SVM margin maximization resists overfitting in high-dimensional sparse text vectors. & Requires manual feature engineering; linear kernel ignores word order. \\
\addlinespace
McCallum \& Nigam (1998) \cite{mccallum1998comparison} & Multinomial Naive Bayes & Established MNB superiority over Bernoulli model; incorporated Laplace smoothing. & Independent feature assumption fundamentally violates natural language grammar. \\
\addlinespace
Go et al. (2009) \cite{go2009twitter} & Distant Supervision Twitter NLP & Automated corpus creation using emoticons; achieved 82.2\% SVM accuracy on 1.6M tweets. & Severe label noise from sarcastic emoticon usage; weak contextual disambiguation. \\
\addlinespace
Hutto \& Gilbert (2014) \cite{hutto2014vader} & Rule-based Lexicon (VADER) & Built specialized social media sentiment lexicon with 5 syntactic heuristic rules. & Static dictionary; misclassifies gaming jargon (``sick'', ``nasty'', ``diff''). \\
\addlinespace
Mohammad \& Turney (2013) \cite{mohammad2013crowdsourcing} & Emotion Lexicon Mining & Crowdsourced word-emotion associations across 8 primary affect categories. & Unsupervised word-level matching cannot resolve context-dependent valence flips. \\
\addlinespace
Rosenthal et al. (2017) \cite{rosenthal2017semeval} & SemEval Twitter Benchmarks & Standardized evaluation protocols for computer-mediated communication sentiment. & Benchmarks concentrated on generic Twitter discourse, omitting gaming jargon. \\
\addlinespace
Vaswani et al. (2017) \cite{vaswani2017attention} & Transformer Architecture & Replaced recurrence with multi-head self-attention; enabled massive parallel pre-training. & Quadratic computational complexity $\mathcal{O}(n^2)$ with respect to sequence length $n$. \\
\addlinespace
Devlin et al. (2019) \cite{devlin2019bert} & Bidirectional Pre-training (BERT) & Introduced masked language modeling; established state-of-the-art across 11 NLP tasks. & 110M parameters; high memory and inference latency (>200\,ms on CPU). \\
\addlinespace
Sanh et al. (2019) \cite{sanh2019distilbert} & Knowledge Distillation & Compressed BERT by 40\% via distillation; retained 97\% language capabilities, 60\% faster. & Still incurs 40--50\,ms CPU latency, too slow for unbatched microsecond web loops. \\
\addlinespace
Liu et al. (2019) \cite{liu2019roberta} & Optimized Pre-training (RoBERTa) & Demonstrated BERT undertraining; improved performance via dynamic masking. & Substantial parameter overhead; computationally intractable for lightweight hosting. \\
\addlinespace
Mora-Cantallops \& Sicilia (2018) \cite{mora2018exploring} & MOBA Player Sentiment & Discovered empirical link between patch changes, champion balance, and forum toxicity. & Relied on manual qualitative sampling and basic keyword frequency counts. \\
\addlinespace
Makarov et al. (2018) \cite{makarov2018predicting} & Esports Match Forecasting & Modeled Dota 2 and CS:GO match outcomes using in-game telemetries and draft data. & Exclusively quantitative box scores; completely lacked natural language sentiment. \\
\addlinespace
Hodge et al. (2019) \cite{hodge2019toxicity} & Competitive FPS Toxicity & Categorized tactical vs. toxic communications in team shooters (CS:GO, Overwatch). & Qualitative survey methodology; lacked an automated production classification engine. \\
\addlinespace
Sun et al. (2020) \cite{sun2020mobilebert} & Multi-Tier Model Cascades & Proved adaptive routing reduces system inference cost by >70\% without accuracy loss. & Investigated generic mobile workloads; did not implement esports domain adaptation. \\
\bottomrule
\end{tabularx}
\end{table}

\section{Gap Analysis}
\label{sec:gap_analysis}

A critical synthesis of the literature survey reveals five foundational technical and architectural gaps:
\begin{enumerate}
    \item \textbf{Absence of Specialized Gaming Lexicons in Commercial NLP:} Mainstream sentiment analysis tools and cloud APIs (Google Cloud NLP, AWS Comprehend) are trained on formal news, product reviews, or standard Twitter corpora. They completely lack semantic representations for tactical esports neologisms (\textit{e.g.}, ``diff'', ``clutch'', ``nerf'', ``whiff'', ``instalock'', ``lineup''), resulting in systemic false-negative and false-positive classifications.
    \item \textbf{Total Disconnect Between Esports Trackers and Sentiment Mining:} Leading esports statistics portals (VLR.gg, Tracker.gg) concentrate exclusively on numerical in-game telemetries (KDA, ACS, headshot percentage). They provide zero semantic understanding of community morale, patch frustration, or qualitative player perception.
    \item \textbf{The Latency vs. Accuracy Trade-Off Dilemma:} Existing sentiment deployments employ monolithic architectures: either deploying lightweight classical models (Naive Bayes, SVM) that are fast but linguistically brittle, or deploying heavy transformers (BERT, RoBERTa) that are accurate but computationally and financially prohibitive for live web serving. No existing gaming platform implements an adaptive multi-tier cascade.
    \item \textbf{The Black-Box Interpretability Barrier:} Most deep learning sentiment models output raw class probabilities without explaining which specific lexical tokens drove the decision. In gaming moderation, human administrators require transparent Explainable AI (XAI) feature attribution to substantiate disciplinary actions or balance adjustments.
    \item \textbf{Lack of Integrated Production Platforms with Enterprise Security:} Existing academic prototypes remain isolated Jupyter notebooks or command-line scripts lacking secure session management, dual-factor authentication (OAuth 2.0 and OTP), and dual-layer database persistence.
\end{enumerate}

\section{Problem Statement}
\label{sec:problem_statement}

Formally, the computational problem addressed by this project can be stated as follows:

\vspace{0.3cm}
\begin{tcolorbox}[colback=red!3!white,colframe=krmured,arc=2mm,title=\textbf{Formal Problem Statement}]
\textit{Given a continuous, high-velocity stream of uncurated, highly informal gaming reviews and social commentary $D = \{d_1, d_2, \dots, d_N\}$ relating to the tactical first-person shooter \textit{Valorant}, where each document $d_i$ contains domain-specific slang, sarcastic valence inversions, and tactical entity mentions (Agents $A = \{a_1, \dots, a_{24}\}$ and Maps $M = \{m_1, \dots, m_{13}\}$):}

\vspace{0.2cm}
\textit{Design, implement, and benchmark an integrated, publication-grade software intelligence platform that:}
\begin{enumerate}
    \item \textit{Accurately classifies document sentiment polarity $y_i \in \{\text{Positive}, \text{Neutral}, \text{Negative}\}$ with macro F1-score exceeding 90\%;}
    \item \textit{Binds mean inference latency to sub-10 milliseconds for standard operational queries;}
    \item \textit{Extracts and correlates granular sentiment distributions across tactical Agents and Maps;}
    \item \textit{Generates real-time, token-level Explainable AI (XAI) visual saliency attributions; and}
    \item \textit{Enforces zero-trust enterprise access via Google OAuth 2.0 and Gmail OTP two-factor verification, backed by synchronized local and cloud database persistence.}
\end{enumerate}
\end{tcolorbox}

\section{Research and Technical Objectives}
\label{sec:objectives}

To resolve the problem statement and bridge the identified research gaps, this project defines six measurable, sequential technical objectives:

\begin{enumerate}
    \item \textbf{Objective 1: Dataset Engineering and Domain Preprocessing Pipeline}
    \begin{itemize}
        \item Aggregate, clean, and validate a stratified corpus of over 100,000 authentic \textit{Valorant} player reviews from Steam Community and Reddit.
        \item Engineer a specialized text normalization pipeline in \texttt{src/text\_processing.py} featuring regex noise reduction, contraction resolution, and valence-aware slang translation.
    \end{itemize}

    \item \textbf{Objective 2: Multi-Tier Hybrid Inference Cascade Design}
    \begin{itemize}
        \item Formulate and train an adaptive 4-tier model hierarchy:
        \begin{itemize}
            \item \textbf{Tier 1:} Sub-linear TF-IDF vectorizer + Platt-calibrated Linear Support Vector Classifier (LinearSVC).
            \item \textbf{Tier 2:} Multinomial Naive Bayes probabilistic model with Laplace smoothing.
            \item \textbf{Tier 3:} Fine-tuned \textit{DistilBERT-base-uncased} contextual transformer trained via PyTorch and CrossEntropyLoss.
            \item \textbf{Tier 4:} Zero-shot Large Language Model (LLM) reasoning fallback for complex idiomatic ambiguities.
        \end{itemize}
    \end{itemize}

    \item \textbf{Objective 3: Entity-Level Sentiment Slicing and Tactical Correlation}
    \begin{itemize}
        \item Correlate classified polarities with in-game tactical entities across 24 playable Agents and 13 competitive Maps to compute community satisfaction indices and patch reception metrics.
    \end{itemize}

    \item \textbf{Objective 4: Explainable AI (XAI) Token Attribution Engine}
    \begin{itemize}
        \item Implement token-level feature attribution algorithms that project model weight vectors and gradient saliencies back to input tokens, providing real-time visual highlight glows (emerald-green for positive, crimson for negative) inside the user interface.
    \end{itemize}

    \item \textbf{Objective 5: Full-Stack Web Application and Enterprise Authentication}
    \begin{itemize}
        \item Develop a responsive, cyberpunk-inspired Tactical HUD web application utilizing Flask, Jinja2, and asynchronous AJAX endpoints.
        \item Implement dual-factor zero-trust authentication integrating Google OAuth 2.0 alongside SMTP-backed Gmail OTP verification with cryptographically signed session tokens.
    \end{itemize}

    \item \textbf{Objective 6: Dual Persistence Architecture and Cloud CI/CD Deployment}
    \begin{itemize}
        \item Engineer a resilient dual-database persistence engine combining local SQLite via SQLAlchemy ORM with cloud PostgreSQL via Supabase PostgREST.
        \item Containerize the application using Docker and deploy to Render Cloud PaaS with automated health checks and zero-cost hosting parameters.
    \end{itemize}
\end{enumerate}
